% -------------------------------------------------------------
% SCRUM IMPLEMENTATION
% -------------------------------------------------------------
\section{Scrum Implementation}

\hspace*{1cm}The team implemented SCRUM and its ceremonies from the very first week of development. Two week long sprints were used for the whole duration of development. However, there were certain changes throughout the semester based on team's experience in order to optimize team's workflow.

\medskip

\hspace*{1cm}Scrum master was elected in the beginning and due to his effectivity he was in this role until the end of this project. All team members agreed on using Jira. Jira was chosen because it allowed the team to use all features such as task management and proper story point assignment without the need of additional subscriptions. Jira proved to be a good fit, because it is better optimized for bigger projects the team found it very helpful. Backlog was used for storing future tasks, from backlog certain tasks were taken to sprint backlog during sprint plannings. Board for each sprint consisted of several columns, “TO DO”, “IN PROGRESS”, “UNIT TESTING”, “UNDER REVIEW” and “DONE”. One member connected Jira seamlessly with IDEs, so when a team member started on a new branch the task would automatically move to “IN PROGRESS”. Very important aspect of team cooperation was the approach to reviews. Every task required a minimum of two reviewers. If a reviewer found something imperfect, a change request was created in GitHub and the branch owner had to fix it before the reviewer would approve it again.

\medskip

\hspace*{1cm}Project was finished in 6 sprints. Each sprint consisted of necessary ceremonies. Our team decided to have two standup meetings, one in each week of the two week sprint, usually on Mondays and Wednesdays. Standup meetings were short, structured and valuable part of our communication throughout the semester and they stayed unchanged until the end.

\medskip
\medskip

\hspace*{1cm}Sprint plannings were held every other week before the start of a new sprint. Four hour block was determined to this ceremony every single time. For the first two sprints plannings were held on Mondays. After this period team decided to move sprint planning to Friday the week before. Monday was particularly busy day for all team members and moving sprint planning to Friday had the benefit of one more weekend, so it essentially prolonged the sprint period. With each new sprint planning less time was needed as the team members understood and got better in planning, focusing and critically thinking about what needs to be done and who is the correct person for specific task.

\medskip

\hspace*{1cm}At the end of each sprint team held two last ceremonies. Sprint review and sprint retrospective. Both of these ceremonies were held on Fridays. 

\hspace*{1cm}Originally, all deadlines for tasks and pull-request reviews were set for Thursday evening. However, the team found this really ineffective and there was not enough time for reviews to be done. The deadline for tasks was later moved to Wednesdays evening and reviews deadline to Thursdays evening, so there would be enough time for all ceremonies held on Friday. 

\hspace*{1cm}In sprint review team members stated what each of them accomplished and how they felt about it. This ceremony was especially important for the team as it provided information about what each member prefers, what is their level at certain tasks and it allowed everyone to understand what exactly and how things were done in the past two weeks. Sprint retrospective was done through Jira. For each sprint one team member created a table with three columns “Start doing”, “Stop doing” and “Keep doing”. After ten minutes of anonymous answering team members checked and talked about the topics in the table. This turned out to be as helpful as sprint review. Retrospective was definitely the ceremony that contributed the most to team's growth and ability to cooperate effectively.
